% 数学ノート集・共通テンプレート v2（積分論2026年度第1部の版面を参考）
% A5判。章内で定義・定理・例・注意等の番号を共通化。
% コンパイル： platex -interaction=nonstopmode math_note_template_v2.tex (2回)
%              dvipdfmx math_note_template_v2.dvi
\documentclass[a5paper,10pt,onecolumn]{jsarticle}
\usepackage[dvipdfmx,top=22mm,bottom=20mm,inner=19mm,outer=18mm,headheight=13pt,headsep=6mm]{geometry}
\usepackage{amsmath,amssymb,amsthm,mathtools}
\usepackage{enumitem,fancyhdr,titlesec}
\usepackage[dvipdfmx,bookmarks=true,colorlinks=false]{hyperref}
\usepackage{pxjahyper}

% ----- 差し替える文書情報 -----
\newcommand{\NoteTitle}{算術級数中の素数とDirichletの定理}
\newcommand{\NoteSeries}{数学ノート集}

% ----- 文字・組版 -----
% 参考PDFの質素な明朝＋Computer Modernの印象を引き継ぎながら，
% A5判で窮屈にならないよう本文の行間を少し確保する．
\linespread{1.13}
\setlength{\parindent}{1zw}
\setlength{\parskip}{0pt}
\setlength{\abovedisplayskip}{6pt plus 2pt minus 2pt}
\setlength{\belowdisplayskip}{6pt plus 2pt minus 2pt}
\setlength{\abovedisplayshortskip}{3pt plus 2pt minus 1pt}
\setlength{\belowdisplayshortskip}{4pt plus 2pt minus 1pt}
\setlist[enumerate]{topsep=3pt,itemsep=1pt,parsep=0pt,leftmargin=2.4zw}
\titleformat{\section}{\normalfont\normalsize\bfseries}{\thesection}{0.8zw}{}
\titlespacing*{\section}{0pt}{2.0ex plus .4ex}{.9ex}
\titleformat{\subsection}{\normalfont\normalsize\bfseries}{\thesubsection}{0.7zw}{}
\titlespacing*{\subsection}{0pt}{1.5ex plus .3ex}{.6ex}
\raggedbottom

% ----- 柱とページ番号：ページ上部、罫線なし -----
\fancypagestyle{mathnote}{%
  \fancyhf{}%
  \fancyhead[L]{\footnotesize \NoteSeries\ （\NoteTitle）}%
  \fancyhead[R]{\footnotesize \thepage}%
  \renewcommand{\headrulewidth}{0pt}%
  \renewcommand{\footrulewidth}{0pt}%
}
\pagestyle{mathnote}

% ----- 節内通し番号の定義・定理群 -----
\newtheoremstyle{mathnotetheorem}
 {4.5pt}{4.5pt}{\normalfont}{}%
 {\bfseries}{}{0.65zw}%
 {\thmname{#1}\thmnumber{ #2}\thmnote{（#3）}}
\theoremstyle{mathnotetheorem}
\newtheorem{theorem}{定理}[section]
\newtheorem{definition}[theorem]{定義}
\newtheorem{lemma}[theorem]{補題}
\newtheorem{proposition}[theorem]{命題}
\newtheorem{corollary}[theorem]{系}
\newtheorem{example}[theorem]{例}
\newtheorem{remark}[theorem]{注意}
\renewcommand{\proofname}{証明}
\renewcommand{\qedsymbol}{$\square$}
\makeatletter
\renewenvironment{proof}[1][\proofname]{%
  \par\pushQED{\qed}\normalfont\topsep4\p@\@plus2\p@\relax
  \trivlist\item[\hskip\labelsep\bfseries #1]\ignorespaces
}{\popQED\endtrivlist\@endpefalse}
\makeatother

% ----- 使うノートで適宜拡張する記号 -----
\newcommand{\R}{\mathbb{R}}
\newcommand{\norm}[1]{\lVert #1\rVert}
\newcommand{\abs}[1]{\lvert #1\rvert}

\begin{document}
\section*{読了後にできること}
本稿を読むための前提は，Dirichlet指標，合同式，Dirichletの定理．
本稿の具体的な到達課題は，Dirichletの定理を具体的な合同類へ適用し，互いに素条件を検査する．

\section{有限Abel群の指標}

本稿では整数 $q\ge2$ を固定する．
$G=(\mathbb Z/q\mathbb Z)^\times$ と置き，
その位数を $\varphi(q)$ と書く．
有限Abel群の構造定理と複素解析の基本定理を前提とする．

\begin{definition}
$G$ から複素数の単位円への群準同型を指標という．
$G$ の指標を，$\gcd(n,q)=1$ なら $[n]\in G$ によって，
$\gcd(n,q)>1$ なら零として整数全体へ延長したものを
\emph{$q$ を法とするDirichlet指標}という．
全体を $\widehat G$ と書く．
$\gcd(n,q)=1$ なら $\chi(n)=1$，それ以外なら零となる指標を
主指標 $\chi_0$ と書く．
\end{definition}

\begin{lemma}[指標の直交関係]\label{lem:orth}
有限Abel群 $G$ について $\#\widehat G=\#G$ であり，
\[
 \sum_{\chi\in\widehat G}\chi(g)
 =\begin{cases}\#G&(g=e),\\0&(g\ne e),\end{cases}
\]
\[
 \sum_{g\in G}\chi(g)\overline{\psi(g)}
 =\begin{cases}\#G&(\chi=\psi),\\0&(\chi\ne\psi).\end{cases}
\]
\end{lemma}
\begin{proof}
有限Abel群の構造定理から
$G\cong C_{n_1}\times\cdots\times C_{n_r}$ と書ける．
$C_n$ の指標は
$[j]\mapsto\exp(2\pi ikj/n)$（$0\le k<n$）の
ちょうど $n$ 個である．
従って直積の指標は成分ごとの指標の積として全て得られ，
個数は $\prod_i n_i=\#G$ である．
$g\ne e$ ならある成分が非零なので，
その成分に関する単位根の和が零となり
最初の直交式を得る．
二番目の式は $\chi\overline\psi$ が指標であり，
非自明ならその群全体での和が零になることから従う．
\end{proof}

\begin{corollary}[合同類の指示関数]\label{cor:indicator}
$\gcd(a,q)=1$ とする．任意の整数 $n$ について
\[
 1_{\{n\equiv a\pmod q\}}
 =\frac1{\varphi(q)}
 \sum_{\chi\bmod q}\overline{\chi(a)}\chi(n).
\]
\end{corollary}
\begin{proof}
$\gcd(n,q)>1$ なら右辺の各 $\chi(n)$ は零であり，
$n\equiv a\pmod q$ でもないので成立する．
$\gcd(n,q)=1$ なら右辺は
$\varphi(q)^{-1}\sum_\chi\chi([n][a]^{-1})$
に等しい．
補題\ref{lem:orth}からこれは $[n]=[a]$ のとき1，
そうでないとき0である．
\end{proof}

\section{Dirichletの $L$ 関数}

\begin{definition}
Dirichlet指標 $\chi$ に対して
\[
 L(s,\chi)=\sum_{n=1}^\infty\frac{\chi(n)}{n^s},
 \qquad \operatorname{Re}s>1
\]
をDirichletの $L$ 関数という．
\end{definition}

\begin{lemma}[非主指標の部分和]\label{lem:partialsums}
$\chi\ne\chi_0$ なら
\[
 \sum_{a=1}^q\chi(a)=0,\qquad
 \sup_{N\ge1}\left|\sum_{n=1}^N\chi(n)\right|\le q.
\]
\end{lemma}
\begin{proof}
$\chi$ は周期 $q$ であり，
一周期の和は
$\sum_{g\in G}\chi(g)$ に等しい．
$\chi$ は非自明なので補題\ref{lem:orth}から零である．
$N=kq+r$，$0\le r<q$ と分けると
$N$ までの和は完全な $k$ 周期の和が零なので
$\sum_{n=1}^r\chi(n)$ に等しい．
$|\chi(n)|\le1$ だから絶対値は $r\le q$ 以下である．
\end{proof}

\begin{theorem}[非主 $L$ 関数の解析接続]\label{thm:L-hol}
$\chi$ が非主指標なら，
$L(s,\chi)$ は $\operatorname{Re}s>0$ で
局所一様収束する級数によって正則関数に延長される．
特に
\[
 L(1,\chi)=\sum_{n=1}^\infty\frac{\chi(n)}n
\]
は通常の意味で収束する．
\end{theorem}
\begin{proof}
$A(N)=\sum_{n\le N}\chi(n)$ と置く．
補題\ref{lem:partialsums}によって $|A(N)|\le q$ である．
Abelの部分和公式から
\[
 \sum_{n=M}^{N}\frac{\chi(n)}{n^s}
 =A(N)N^{-s}-A(M-1)M^{-s}
  +s\int_M^N A(\lfloor x\rfloor)x^{-s-1}\,dx
\]
となる．
$\operatorname{Re}s\ge\varepsilon>0$，
$|s|\le R$ のコンパクト集合では右辺の尾部は
$O_{\varepsilon,R,q}(M^{-\varepsilon})$ で抑えられる．
従って級数は $\operatorname{Re}s>0$ で局所一様収束し，
各項が正則であることから和も正則である．
$s=1$ においても同じ収束が成立する．
\end{proof}

\begin{proposition}[Euler積]\label{prop:euler}
$\operatorname{Re}s>1$ なら
\[
 L(s,\chi)=\prod_{p\ {\rm prime}}
 (1-\chi(p)p^{-s})^{-1}.
\]
特にこの半平面では $L(s,\chi)\ne0$ である．
\end{proposition}
\begin{proof}
有限個の素数について幾何級数を展開すると，
完全乗法性によって，それらの素数だけを素因数とする
$n$ の項 $\chi(n)n^{-s}$ の和を得る．
$\sum_n|\chi(n)|n^{-\sigma}\le\zeta(\sigma)<\infty$
（$\sigma=\operatorname{Re}s>1$）なので，
有限素数集合を増加させると絶対収束によって
Dirichlet級数へ収束する．
さらに
\[
 \sum_p\sum_{k\ge1}
 \frac{|\chi(p)|^k}{k p^{k\sigma}}
 \le\sum_p\frac{p^{-\sigma}}{1-p^{-\sigma}}<\infty
\]
だから対数Euler積も絶対収束し，
その指数は零でない．
\end{proof}

\section{主指標とRiemannゼータ関数}

\begin{proposition}[ゼータ関数の局所的な解析接続]
\label{prop:zeta}
$\operatorname{Re}s>0$ において
\[
 \zeta(s)=\frac{s}{s-1}
 -s\int_1^\infty\{x\}x^{-s-1}\,dx
 \tag{2}
\]
は，$s=1$ に留数1の単純極を持つ正則関数の有理型延長を与える．
また
\[
 L(s,\chi_0)=\zeta(s)\prod_{p\mid q}(1-p^{-s})
\]
である．
従って $L(s,\chi_0)$ は $s=1$ に単純極を持ち，
その留数は $\prod_{p\mid q}(1-p^{-1})>0$ である．
\end{proposition}
\begin{proof}
$\operatorname{Re}s>1$ では部分積分によって
\[
 \zeta(s)=s\int_1^\infty\lfloor x\rfloor x^{-s-1}\,dx
\]
が成立する．
$\lfloor x\rfloor=x-\{x\}$ を用いれば式 (2) を得る．
$0\le\{x\}<1$ なので積分は
$\operatorname{Re}s>0$ のコンパクト集合上で局所一様収束し，
正則関数を定める．
第一項 $s/(s-1)$ は $s=1$ に留数1の単純極を持つ．
主指標のEuler積から，$q$ を割る素数についての
Euler因子を取り除けば最後の積表示を得る．
有限積は $s=1$ で非零なので，
留数に関する主張が従う．
\end{proof}

\section{非零性の準備：正係数Dirichlet級数}

\begin{lemma}[Landauの正係数補題]\label{lem:landau}
$a_n\ge0$ とし
$F(s)=\sum_{n\ge1}a_n n^{-s}$ の収束半平面が
$\operatorname{Re}s>\sigma_c$ で，
$\sigma_c$ は有限とする．
このとき $F$ は実点 $s=\sigma_c$ の近傍へ
正則に延長できない．
\end{lemma}
\begin{proof}
背理法で，$F$ が $\sigma_c$ の近傍まで正則に延長するとする．
$\sigma_0>\sigma_c$ を十分近く取れば，
$\sigma_0$ を中心とするTaylor級数は
$\sigma_0-\delta<\sigma_c$ となる正数 $\delta$ まで
収束するように選べる．
収束半平面の内部では級数を項別微分でき，
\[
 F^{(k)}(\sigma_0)
 =\sum_{n\ge1}a_n(-\log n)^k n^{-\sigma_0}.
\]
従ってTaylor級数を $\sigma_0-\delta$ で評価し，
各項が非負であることに注意してTonelliの定理を適用すると
\[
 \begin{aligned}
 F(\sigma_0-\delta)
 &=\sum_{k=0}^\infty\frac{\delta^k}{k!}
    \sum_{n\ge1}a_n(\log n)^k n^{-\sigma_0}\\
 &=\sum_{n\ge1}a_n n^{-\sigma_0}
    \sum_{k=0}^\infty\frac{(\delta\log n)^k}{k!}\\
 &=\sum_{n\ge1}a_n n^{-(\sigma_0-\delta)}.
 \end{aligned}
\]
左辺は有限なので右辺も収束する．
これは $\sigma_0-\delta<\sigma_c$ に矛盾する．
\end{proof}

\begin{lemma}[実非主指標の非零性]\label{lem:real-nonzero}
$\chi$ が実数値の非主Dirichlet指標なら
$L(1,\chi)>0$ である．
\end{lemma}
\begin{proof}
定理\ref{thm:L-hol}から $L(1,\chi)$ は実数である．
$\sigma>1$ に対して
\[
 \zeta(\sigma)L(\sigma,\chi)
 =\sum_{n\ge1}\frac{a_n}{n^\sigma},
 \qquad
 a_n=\sum_{d\mid n}\chi(d)
\]
と書ける．
$a_n$ は乗法的であり，
素数冪 $p^k$ について
$a_{p^k}=\sum_{j=0}^k\chi(p)^j$ である．
$\chi(p)$ は $0,1,-1$ のいずれかなので
$a_{p^k}\ge0$ である．
また指数 $2k$ が偶数なら
$a_{p^{2k}}\ge1$ となる．
従って全ての $n$ について $a_n\ge0$ かつ
$a_{n^2}\ge1$ である．
このため
$\sum_n a_nn^{-1/2}\ge\sum_m m^{-1}=\infty$ となり，
このDirichlet級数の収束境界 $\sigma_c$ は
少なくとも $1/2$ である．

もし $L(1,\chi)=0$ なら，
命題\ref{prop:zeta}によって
$\zeta(s)L(s,\chi)$ の $s=1$ の極が相殺され，
この積は半平面 $\operatorname{Re}s>0$ で正則となる．
級数の収束境界は $\sigma_c\le1$ なので，
特に実点 $\sigma_c\in[1/2,1]$ においても
正則である．これは補題\ref{lem:landau}に矛盾する．
従って $L(1,\chi)\ne0$ である．
$\sigma>1$ ではEuler積から
$L(\sigma,\chi)>0$ であり，
$\sigma\downarrow1$ の極限も非負なので
$L(1,\chi)>0$ となる．
\end{proof}

\section{非実指標の非零性}

\begin{lemma}[三角多項式による非零性]\label{lem:complex-nonzero}
$\chi$ を非実Dirichlet指標とする．
このとき $L(1,\chi)\ne0$ である．
\end{lemma}
\begin{proof}
$\chi$ が非実なので $\chi^2\ne\chi_0$ である．
実際 $\chi^2=\chi_0$ なら
$\gcd(n,q)=1$ の元に対し $\chi(n)\in\{\pm1\}$ となり，
$\chi$ は実数値となる．
$\sigma>1$ に対しEuler積の対数を用いると
\[
 \begin{aligned}
 &\log\Bigl(
 \zeta(\sigma)^3|L(\sigma,\chi)|^4
 |L(\sigma,\chi^2)|\Bigr)\\
 &=\sum_p\sum_{k=1}^\infty
 \frac{3+4\operatorname{Re}(\chi(p)^k)
   +\operatorname{Re}(\chi(p)^{2k})}
      {k p^{k\sigma}}
 \end{aligned}
\]
となる．
$\chi(p)=0$ なら分子は3であり，
$|\chi(p)|=1$ なら $\chi(p)^k=e^{i\theta}$ と書いて
\[
 3+4\cos\theta+\cos2\theta
 =2(1+\cos\theta)^2\ge0.
\]
従って対数は非負であり，
\[
 \zeta(\sigma)^3|L(\sigma,\chi)|^4
 |L(\sigma,\chi^2)|\ge1.
 \tag{3}
\]

もし $L(1,\chi)=0$ なら
正則性により $L(\sigma,\chi)=O(\sigma-1)$ である．
$\chi^2$ は非主指標なので
$L(\sigma,\chi^2)$ は $\sigma\downarrow1$ で有界であり，
$\zeta(\sigma)=O((\sigma-1)^{-1})$ である．
従って式 (3) の左辺は
$O((\sigma-1)^{-3+4})=O(\sigma-1)$ となり，
零へ収束する．これは左辺が常に1以上という不等式に反する．
ゆえに $L(1,\chi)\ne0$ である．
\end{proof}

\begin{theorem}[非主指標の非零性]\label{thm:nonvanishing}
$q$ を法とする任意の非主Dirichlet指標 $\chi$ に対し
\[
 L(1,\chi)\ne0.
\]
\end{theorem}
\begin{proof}
$\chi$ が実数値なら補題\ref{lem:real-nonzero}，
非実なら補題\ref{lem:complex-nonzero}を適用すればよい．
\end{proof}

\section{素数の対数Euler積}

\begin{lemma}[対数Euler積の主項]\label{lem:prime-log}
$\sigma\downarrow1$ のとき
\[
 \sum_p\frac{\chi_0(p)}{p^\sigma}
 =\log\frac1{\sigma-1}+O_q(1)
\]
であり，非主指標 $\chi$ に対して
\[
 \sum_p\frac{\chi(p)}{p^\sigma}=O_{q,\chi}(1)
\]
が成り立つ．
\end{lemma}
\begin{proof}
$\sigma>1$ において
\[
 \log L(\sigma,\chi)
 =\sum_p\frac{\chi(p)}{p^\sigma}
  +\sum_p\sum_{k\ge2}\frac{\chi(p)^k}{k p^{k\sigma}}
\]
が成立する．
第二項の絶対値は $\sigma\ge1$ で
\[
 \sum_p\sum_{k\ge2}\frac1{k p^k}
 \le\sum_{n=2}^\infty\frac1{n(n-1)}
 <\infty
\]
で一様に抑えられる．
主指標については
命題\ref{prop:zeta}から
$L(\sigma,\chi_0)\sim c_q/(\sigma-1)$，
$c_q>0$ なので
$\log L(\sigma,\chi_0)
 =\log(1/(\sigma-1))+O_q(1)$ となる．
非主指標では定理\ref{thm:nonvanishing}から
$L(\sigma,\chi)\to L(1,\chi)\ne0$ であり，
その近傍で対数の連続な枝を取れるため
$\log L(\sigma,\chi)=O_{q,\chi}(1)$ である．
従って主張を得る．
\end{proof}

\begin{theorem}[Dirichletの算術級数定理]\label{thm:dirichlet}
$\gcd(a,q)=1$ なら，合同条件
\[
 p\equiv a\pmod q
\]
を満たす素数 $p$ は無限に存在する．
さらに
\[
 \sum_{\substack{p\ {\rm prime}\\p\equiv a\pmod q}}
 p^{-\sigma}
 =\frac1{\varphi(q)}
 \log\frac1{\sigma-1}+O_{a,q}(1)
 \quad(\sigma\downarrow1)
\]
が成り立つ．
\end{theorem}
\begin{proof}
系\ref{cor:indicator}を素数に適用し，
$\sigma>1$ で素数の総和を取ると
\[
 \sum_{p\equiv a\pmod q}p^{-\sigma}
 =\frac1{\varphi(q)}
 \sum_{\chi\bmod q}\overline{\chi(a)}
 \sum_p\chi(p)p^{-\sigma}.
\]
指標は有限個なので
補題\ref{lem:prime-log}を各項に適用できる．
主指標では $\overline{\chi_0(a)}=1$ であり，
非主指標の寄与は全て有界である．
従って表示した漸近式を得る．
右辺は $\sigma\downarrow1$ で無限大に発散する．
もし該当する素数が有限個しかなければ，
左辺は有限個の連続な関数の和であり有界となる．
矛盾するので素数は無限に存在する．
\end{proof}

\begin{example}
$4$ を法として $1$ および $3$ と合同な素数は
いずれも無限に存在する．
\end{example}
\begin{proof}
$1,3$ はいずれも $4$ と互いに素なので
定理\ref{thm:dirichlet}を適用できる．
\end{proof}

\begin{remark}
この証明で本質的なのは
非主指標の $L(1,\chi)\ne0$ である．
非実指標では三角多項式の非負性，
実指標では正係数Dirichlet級数のLandau補題を用い，
二種類の証明を区別した．
本定理は各合同類に素数が無限にあることを示すが，
素数の個数の精密な漸近式は別の定理を要する．
\end{remark}

\section{到達課題の完全解答}
\noindent\textbf{算術級数中の素数の判定条件．}
\begin{proposition}
$4n+1$ の形の素数と $4n+3$ の形の素数は
それぞれ無限個存在する．
\end{proposition}
\begin{proof}
$\gcd(1,4)=\gcd(3,4)=1$ である．
Dirichletの定理をそれぞれ法4の剰余類
$1$，$3$ に適用すれば，
どちらの算術級数にも無限個の素数が存在する．
\end{proof}
\begin{proposition}
$6n+2$ の形の素数は2だけである．
\end{proposition}
\begin{proof}
$6n+2=2(3n+1)$ であるから，
$n\ge1$ では二つの正整数の積となり合成数である．
$n=0$ で2となる．
\end{proof}
\noindent\textbf{到達確認．}
$10n+3$ の形の素数は無限に存在するか．
\par\noindent\textbf{解答．}
$\gcd(10,3)=1$ だからDirichletの定理から
無限に存在する．定理を適用する際，
法と剰余が互いに素であることの確認が必要である．
\end{document}